% ============================================================
\chapter{Interior Point Methods}
\label{ch:interior_point}
% ============================================================

\section{Introduction}

In 1984, Narendra Karmarkar, an engineer at AT\&T Bell Labs, published an algorithm that shook the optimization world: a polynomial-time linear programming method that, unlike Dantzig's simplex, does not traverse the vertices of the polyhedron but passes through its \emph{interior}. The community was electrified --- and divided. Some saw a revolution; others dismissed it as hype (the simplex, though exponential in the worst case, is fast in practice). But history would decide: interior point methods, refined by Nesterov and Nemirovski in the 1990s, proved decisive for semidefinite optimization, conic optimization, and large-scale problems. The idea is elegant: replace inequality constraints with a logarithmic barrier, then follow the \emph{central path} toward the optimum.

\begin{intuition}
Imagine you are inside a room (the feasible polyhedron) searching for the
lowest point (the optimum). The simplex walks along the walls and corners.
Interior point methods walk through the middle of the room, gradually
approaching the optimal solution (often near a wall). An invisible barrier
prevents you from touching the walls, but this barrier weakens over
iterations.
\end{intuition}

\section{Barrier method}

\subsection{Formulation}

\begin{definition}[Inequality-constrained problem]
Consider the convex problem:
\[
  \min_{x \in \R^n} f_0(x) \quad \text{s.t.} \quad
  f_i(x) \leq 0,\; i = 1, \ldots, m, \qquad Ax = b,
\]
where $f_0, f_1, \ldots, f_m$ are convex and twice differentiable.
\end{definition}

\begin{definition}[Logarithmic barrier function]
\index{logarithmic barrier}
The \textbf{logarithmic barrier function} is:
\[
  \phi(x) = -\sum_{i=1}^m \ln(-f_i(x)),
\]
defined on the strict interior $\{x : f_i(x) < 0,\; \forall i\}$.
\end{definition}

\begin{definition}[Barrier problem]
\index{barrier method}
The \textbf{barrier problem} for parameter $t > 0$ is:
\[
  \min_x\; t\, f_0(x) + \phi(x) \quad \text{s.t.} \quad Ax = b.
\]
Its solution $x^*(t)$ is called the \textbf{central point} for parameter $t$.
\end{definition}

\subsection{Central path}

\begin{definition}[Central path]
\index{central path}
The \textbf{central path} is the set $\{x^*(t) : t > 0\}$. As
$t \to +\infty$, $x^*(t) \to x^*$, the solution of the original problem.
\end{definition}

\begin{theorem}[Suboptimality bound]
For any $t > 0$, the central point $x^*(t)$ satisfies:
\[
  f_0(x^*(t)) - p^* \leq \frac{m}{t},
\]
where $p^*$ is the optimal value and $m$ is the number of inequality
constraints.
\end{theorem}

\begin{proof}
The KKT conditions of the barrier problem show that $x^*(t)$ is feasible for
the original problem, with multipliers
$\lambda_i^*(t) = -1/(t\, f_i(x^*(t)))$. By weak duality:
\[
  f_0(x^*(t)) - p^* \leq \sum_{i=1}^m \lambda_i^*(t)(-f_i(x^*(t)))
  = \sum_{i=1}^m \frac{1}{t} = \frac{m}{t}.
\]
\end{proof}

\begin{algorithme}[title=\textbf{Barrier Method}]
\begin{enumerate}
  \item Choose strictly feasible $x_0$, $t_0 > 0$, factor $\mu > 1$,
        tolerance $\varepsilon > 0$.
  \item For $k = 0, 1, 2, \ldots$:
    \begin{enumerate}
      \item \textbf{Centering}: solve (via Newton)
            $\min_x t_k f_0(x) + \phi(x)$ s.t.\ $Ax = b$,
            starting from $x_k$, to obtain $x_{k+1} = x^*(t_k)$.
      \item \textbf{Stopping}: if $m / t_k < \varepsilon$, \textbf{stop}.
      \item \textbf{Update}: $t_{k+1} = \mu \cdot t_k$.
    \end{enumerate}
\end{enumerate}
\end{algorithme}

\begin{keyformulas}[title=\textbf{Barrier method complexity}]
\begin{itemize}
  \item Number of outer iterations (barrier steps):
        $\lceil \frac{\ln(m / (t_0 \varepsilon))}{\ln \mu} \rceil$.
  \item Typical choice: $\mu = 10$ to $20$, yielding few outer iterations.
  \item Each centering step requires a few Newton steps (6--40 in practice).
\end{itemize}
\end{keyformulas}

\section{Short-step and long-step methods}

\begin{definition}[Short-step method]
\index{short-step method}
The \textbf{short-step} method chooses $\mu$ close to 1 (e.g.,
$\mu = 1 + 1/\sqrt{m}$) and performs a single Newton step per centering
phase. The iteration count is $\mathcal{O}(\sqrt{m} \ln(m/\varepsilon))$.
\end{definition}

\begin{definition}[Long-step method]
\index{long-step method}
The \textbf{long-step} method chooses a large $\mu$ and performs multiple
Newton steps to re-center. The total Newton step count is also
$\mathcal{O}(\sqrt{m} \ln(m/\varepsilon))$ but with a larger constant.
\end{definition}

\begin{theorem}[Polynomial complexity]
\index{polynomial complexity}
Interior point methods solve a linear program with $n$ variables and $m$
constraints in $\mathcal{O}(\sqrt{m} \ln(1/\varepsilon))$ Newton steps, each
costing $\mathcal{O}(n^3)$ (or less when exploiting sparsity). The total
complexity is polynomial.
\end{theorem}

\section{Primal-dual interior point method}

\begin{definition}[Modified KKT system]
The \textbf{primal-dual} method simultaneously solves the perturbed KKT
conditions:
\[
  \begin{cases}
    \nabla f_0(x) + \sum_i \lambda_i \nabla f_i(x) + A^\top \nu = 0, \\
    -\lambda_i f_i(x) = 1/t, \quad i = 1, \ldots, m, \\
    Ax = b.
  \end{cases}
\]
Newton's method is applied to this system.
\end{definition}

\begin{remark}
The primal-dual method is more efficient in practice than the pure barrier
method because it does not require solving the centering problem exactly at
each stage. Modern solvers (MOSEK, CVXOPT) use this approach.
\end{remark}

\begin{theorem}[Primal-dual convergence]
The primal-dual method achieves $\varepsilon$-accuracy in
$\mathcal{O}(\sqrt{m} \ln(1/\varepsilon))$ Newton iterations under standard
regularity assumptions.
\end{theorem}

\section{Application to linear programming}

\begin{example}[LP in standard form]
For the LP $\min\; c^\top x$ s.t.\ $Ax = b$, $x \geq 0$, the barrier
function is:
\[
  \phi(x) = -\sum_{j=1}^n \ln(x_j).
\]
The Newton system for centering is:
\[
  \begin{pmatrix}
    X^{-2} & A^\top \\
    A & 0
  \end{pmatrix}
  \begin{pmatrix} \Delta x \\ \Delta \nu \end{pmatrix}
  = -\begin{pmatrix} t\, c - X^{-1} e \\ 0 \end{pmatrix},
\]
where $X = \mathrm{diag}(x)$ and $e = (1, \ldots, 1)^\top$.
\end{example}

\begin{example}[Semidefinite programming]
Interior point methods extend naturally to \textbf{semidefinite programs}
(SDPs), where the constraint $x \geq 0$ is replaced by $X \succeq 0$
(positive semidefinite matrix). The barrier function becomes
$\phi(X) = -\ln \det(X)$.
\end{example}

\section{Exercises}

\begin{exercise}[$\star$ -- Logarithmic barrier]
For the LP $\min -x_1 - x_2$ s.t.\ $x_1 + x_2 \leq 1$, $x_1, x_2 \geq 0$,
write the barrier function and compute $x^*(t)$ for $t = 1, 10, 100$. Verify
that $x^*(t) \to x^*$.
\end{exercise}

\begin{exercise}[$\star\star$ -- Central path]
Plot the central path for $\min x_1$ s.t.\ $x_1^2 + x_2^2 \leq 1$ for
$t \in \{0.1, 1, 10, 100\}$. Show it converges to $(-1, 0)$.
\end{exercise}

\begin{exercise}[$\star\star$ -- Newton step]
Derive the Newton system for the barrier method applied to the constrained
QP: $\min \frac{1}{2} x^\top Q x + c^\top x$ s.t.\ $Ax \leq b$.
\end{exercise}

\begin{exercise}[$\star\star\star$ -- Suboptimality bound]
Prove in detail that $f_0(x^*(t)) - p^* \leq m/t$. Deduce the number of
outer iterations needed to achieve accuracy $\varepsilon$ with factor $\mu$.
\end{exercise}

\begin{exercise}[$\star\star\star$ -- Primal-dual implementation]
Implement the primal-dual interior point method for LP in standard form.
Test on a problem with 50 variables and 20 constraints. Compare iteration
counts with the pure barrier method.
\end{exercise}
